\section{Limitaciones y trabajo futuro}

La evaluación tiene cinco límites principales.

\begin{enumerate}
\item \textbf{Población y etiquetas.} “Natural” describe la prevalencia del split completo, no una muestra probabilística de YouTube. La adquisición incluye búsquedas dirigidas y los rótulos son mayormente asistidos; se requiere una cohorte temporal con doble revisión humana independiente y acuerdo por categoría \cite{artstein2008agreement,moreno2012datasetshift}.
\item \textbf{Incertidumbre del ranking.} El bootstrap agrupa por video, pero no mide variación entre semillas. El empate con el retador impide declarar superioridad estadística. Deben repetirse entrenamientos y evaluarse estabilidad temporal \cite{dror2018significance,cawley2010selection}.
\item \textbf{Salvaguarda permisiva.} El margen de no inferioridad de macro-AUPRC fue 0,10. Aunque estaba congelado antes de test, es amplio para una escala 0--1; una réplica confirmatoria debería fijar un margen sustantivo más estricto.
\item \textbf{Capacidad humana.} La tasa de revisión de 34,8\% satisface el máximo experimental de 40\%, pero todavía no se ha traducido en dotación, duración ni costo. El piloto debe medir tiempo por alerta, falsos negativos residuales y desacuerdo humano.
\item \textbf{Modalidad y deriva.} El sistema observa subtítulos y metadatos, no imagen, audio ni conducta de red. Cambios de canal, plataforma o periodo pueden degradar calibración; se requiere monitoreo de deriva, latencia, privacidad y seguridad \cite{tatman2017captions,aoir2020ethics}.
\end{enumerate}

\section{Conclusiones}

El objetivo general se cumple al construir un moderador semiautomático medible y utilizable que transforma subtítulos en alertas temporales, compara representaciones complementarias y deriva incertidumbre a una persona. Su alcance es priorización y apoyo a la decisión, no sanción autónoma.

El primer objetivo específico se verifica con un corpus v2.1 de 173\,240 chunks, 4\,906 videos y cinco salidas, con 14\,163 chunks dañinos y 2\,709 multietiqueta. La separación por video y los 55\,966 eventos de revisión hacen reproducible la procedencia de cada unidad. Cualitativamente, el proyecto deja de depender de lotes dispersos y dispone de un objeto de estudio versionado.

El segundo objetivo se verifica con 28 individuos y cinco ensembles bajo la misma validation OOF. El promedio suave obtiene BA 0,840, macro-AUPRC 0,555 y macro-F1 0,568; supera puntualmente al mejor individuo en 0,0086 de BA y 0,0392 de AUPRC. La evidencia inferencial permanece inconclusa frente al promedio ponderado, por lo que la conclusión correcta es que la regla lexicográfica selecciona un artefacto reproducible, no que exista un ganador universal.

El tercer objetivo se verifica en la apertura única de 22\,684 chunks de test. La compuerta alcanza BA 0,846 y sensibilidad 0,906. La política deriva 34,8\% y deja 65,2\% automático con BA 0,940; así, la incertidumbre y la carga humana pasan de ser supuestos a cantidades observables. El resultado es favorable para un piloto en modo sombra, aunque la capacidad real debe validarse.

El cuarto objetivo se materializa al integrar la regresión logística ganadora, la cascada E5, Qwen3--LoRA y el ensemble suave con calibradores, umbrales, evidencia temporal y registro de revisión. La variedad de modelos permite contrastar una señal léxica rápida, una arquitectura contextual de dos etapas y un LLM adaptado, mientras el ensemble concentra la decisión recomendada.

En conjunto, el sistema presenta un desempeño competitivo para detección y priorización de contenido peruano: F1 binario 0,656 en la vista 4:1, comparable en magnitud con resultados publicados de tareas relacionadas, y BA 0,846 bajo la prevalencia natural del corpus. Su principal fortaleza es sistémica: datos auditables, comparación común, selección congelada, test separado y supervisión humana conviven en un mismo flujo. El siguiente paso es validar ese flujo prospectivamente con capacidad, costos y deriva reales.

\section*{Disponibilidad y ética}

El código, los cuadernos y la documentación se conservan en \url{https://github.com/lkoc/Trabajo_PLN-MIA-Grupo4}; el corpus, pesos y resultados de mayor tamaño se mantienen en el \href{https://drive.google.com/drive/folders/1H3PKr7CWd-RwavZgEnfLhh94eq-NO9P3}{repositorio de Google Drive con acceso controlado}. Los textos proceden de subtítulos públicos; su redistribución exige revisar términos de plataforma, privacidad, licencias y derechos sobre los subtítulos \cite{youtube2023terms,aoir2020ethics}. Una alerta informa al supervisor y la decisión final permanece bajo control humano.
